% New theorem: predecessor in the one-state successor-seeded recursor.
% Requires the independently audited comparator section sr_comparator.tex.
% All arguments in this section are in ordinary intuitionistic FOL equality
% with numerical coverage, 0 != S0, and contextual literal-equation induction.
\section{Removing the independent successor constructor with one update slot}
\label{sec:sr-predecessor}

The successor-seeded one-state proposal in the shared-continuation draft
left a precise problem unresolved: construct a term $P$ and prove
$P(Sa)=a$ without assuming the ordinary initialized iterator or the
unshifted selector that $P$ is intended to recover.  The following
construction resolves that problem for the ordinary first-order theory.
It uses a finite-cycle update with an explicitly accessible transient
seed, rather than assuming that an arbitrary initializer can be decoded.

\subsection{Exact source theory and the comparator prerequisite}

The numerical terms are
\[
 t::=x\mid0\mid R(n,a;u.h),\qquad Sa:=R(0,a;u.0).
\]
The update binds exactly one numerical variable $u$; all other variables
are fixed passive parameters.  The only computation schemas are
\[
 R(0,a;u.h)=Sa,\qquad
 R(Sn,a;u.h)=h[R(n,a;u.h)/u].
 \tag{SR}\label{eq:sr-source}
\]
Use ordinary intuitionistic first-order equality logic with genuine
falsity, the two boundary sentences
\[
 \forall x\,(x=0\vee\exists y\,x=Sy),\qquad 0=S0\longrightarrow\bot,
\]
and contextual induction on a literal equation:
\[
 \frac{\Gamma\vdash E(0)\qquad\Gamma,E(z)\vdash E(Sz)}
      {\Gamma\vdash E(t)}.
 \tag{EI}\label{eq:sr-ei}
\]
The eigenvariable $z$ is absent from the retained context, the conclusion
term, and all passive parameters.  All descriptions are fixed finite
syntax; an equivalent staged first-order signature gives one symbol
for each earlier update description.  There is no uniform total evaluator.
Denote this ordinary theory by $T_{\rm SR}$.

Write
\[
 T(n,a,b):=R(n,a;u.b),\qquad Z(n):=R(n,0;u.0).
\]
Their equations are
\[
 T(0,a,b)=Sa,\quad T(Sn,a,b)=b,\qquad
 Z(0)=1,\quad Z(Sn)=0.
\]
Here $1=S0$.  Coverage proves $Z(t)=0\vee Z(t)=1$.  Applying $Z$ to
$Sx=0$ gives $0=1$, so successor separation is available without
successor injectivity.

The preceding comparator construction supplies a finite SR term $L$
with whole-domain proofs of
\[
\begin{gathered}
 L(x,y)=0\vee L(x,y)=1,\qquad
 L(0,y)=1,\qquad L(Sx,0)=0,\\
 L(Sx,Sy)=L(x,y),\qquad
 L(x,y)=1\ \longrightarrow\ L(x,Sy)=1.
\end{gathered}
 \tag{L}\label{eq:sr-comparator-properties}
\]
Its proof uses only SR, numerical cases, equality, and literal-equation
induction.  In particular it assumes no predecessor, unshifted selector,
pairing, successor injectivity, or induction on compound formulas.
This dependency is essential: the familiar initialized-scalar
predecessor proof cannot simply be imported before its initializer has
been recovered.

\subsection{Shifted addition and elementary comparison facts}

Define the shifted sum
\[
                  A(a,b):=R(b,a;u.Su).
\]
Its computation equations and elementary literal-equation inductions give
\[
\begin{gathered}
 A(a,0)=Sa,\qquad A(a,Sb)=S A(a,b),\\
 A(0,b)=Sb,\qquad A(Sa,b)=S A(a,b),\qquad A(a,b)=A(b,a).
\end{gathered}
 \tag{A}\label{eq:sr-shifted-addition}
\]
For the third equation induct on $b$; for the fourth do the same with
$a$ fixed.  Commutativity follows by induction on $b$, using those
two equations for its base and step.  These are equations in the source
theory, not identities inferred from the intended value $a+b+1$.
An immediate consequence is
\[
 A(0,1)=2,\qquad
 A(Sa,S^2a)=S^2A(a,Sa).
 \tag{AD}\label{eq:sr-double-index}
\]

Literal-equation induction on $n$ in the comparator recurrences yields
\[
                 L(n,n)=1,\qquad L(Sn,n)=0.
 \tag{Diagonal}\label{eq:sr-comparison-diagonals}
\]
Monotonicity then gives $L(a,Sa)=1$.
Also $L(S^2m,m)=0$: its Boolean alternative $1$ would imply
$L(S^2m,Sm)=1$ by monotonicity, whereas the recurrence and the strict
diagonal make this $L(Sm,m)=0$.  Separation rules out $0=1$.
Finally, monotonicity and the recurrence prove
\[
           L(St,m)=1\ \longrightarrow\ L(t,m)=1,
 \tag{Down}\label{eq:sr-comparison-down}
\]
because $L(St,Sm)=L(t,m)$.  None of these arguments uses successor
cancellation or reflection of the comparator into actual equality.

\subsection{A reset cycle with an accessible transient seed}

For each fixed passive parameter $m$, define
\[
           F_m(u):=T\bigl(Z(L(Su,m)),u,0\bigr).
 \tag{F}\label{eq:sr-cycle-update}
\]
This is a finite SR term.  Its two Boolean cases are
\[
 F_m(u)=Su\quad\hbox{if }L(Su,m)=1,\qquad
 F_m(u)=0\quad\hbox{if }L(Su,m)=0.
\]
Thus it increments below the chosen boundary and resets to zero at or
beyond that boundary.  In particular the preceding comparison facts give
\[
                       F_m(m)=0,\qquad F_m(Sm)=0.
 \tag{Reset}\label{eq:sr-cycle-reset}
\]
The second equality makes the source's available initializer $Sm$
useful: one explicit application of $F_m$ turns it into zero.

Define two derived trajectories
\[
\begin{aligned}
 J_m(t)&:=F_m\bigl(R(t,m;u.F_m(u))\bigr),\\
 K_m(t)&:=F_m\bigl(R(t,m;u.F_m(F_m(u)))\bigr).
\end{aligned}
 \tag{JK}\label{eq:sr-cycle-trajectories}
\]
Their recurrences follow from SR and Reset by direct equational reasoning:
\[
\begin{gathered}
 J_m(0)=0,\qquad J_m(St)=F_m(J_m(t)),\\
 K_m(0)=0,\qquad K_m(St)=F_m(F_m(K_m(t))).
\end{gathered}
 \tag{JR,KR}\label{eq:sr-cycle-recurrences}
\]
For $K$, the outer $F_m$ commutes with the update $F_m\circ F_m$
simply by associativity of substitution.  No general iterator-rotation
theorem or arbitrary initializer has been assumed.  Every induction
below holds $m$ fixed until the completed universal identity is
subsequently specialized.

\subsection{A conditional equation induction, with its literal invariant}

The crucial bounded-trajectory statement is
\[
                   L(t,m)=1\ \longrightarrow\ J_m(t)=t.
 \tag{Run}\label{eq:sr-cycle-bounded-run}
\]
It is proved by induction on the following \emph{single literal equation}:
\[
 T\bigl(L(t,m),0,J_m(t)\bigr)
       =T\bigl(L(t,m),0,t\bigr).
 \tag{Gate}\label{eq:sr-cycle-gate}
\]
When $L(t,m)=0$, both sides are $1$.  When $L(t,m)=1$, the equation
is exactly $J_m(t)=t$.  This use of the shifted selector is valid even
though the ordinary unshifted selector has not yet been constructed.

At $t=0$, the comparator value is $1$ and $J_m(0)=0$, so Gate holds.
For the successor step retain Gate as the literal induction hypothesis
and split the proved Boolean value $L(St,m)$.
If it is $0$, the new gate equation has both sides equal to $1$.
If it is $1$, Down gives $L(t,m)=1$; the retained gate equation
therefore yields $J_m(t)=t$.  The definition of $F_m$ now gives
$F_m(t)=St$, and JR gives $J_m(St)=St$.  This establishes the new
gate equation.  EI proves Gate for every $t$, and its positive Boolean
case proves Run.

The retained induction context contains only the displayed gate
equation and fixed passive data.  Induction on the implication Run,
induction on a disjunction, and changing a parameter inside an active
induction hypothesis are not used.

At the diagonal, Run and $L(m,m)=1$ give $J_m(m)=m$.  One more
computation and Reset consequently give
\[
                          J_m(Sm)=0.
 \tag{Period base}\label{eq:sr-cycle-period-base}
\]

\subsection{Two ordinary equation inductions give the predecessor}

First, with $m$ fixed, induct on the literal equation
\[
                     J_m(A(m,t))=J_m(t).
 \tag{Period}\label{eq:sr-cycle-period}
\]
At zero, $A(m,0)=Sm$ and both trajectories are zero by Period base.
At a successor, $A(m,St)=S A(m,t)$ and JR applies $F_m$ to both
sides of the induction equation.  This proves Period on the entire
source domain; it is not a calculation only at standard counters.

Second, with the same fixed parameter $m$, induct on
\[
                  K_m(Sa)=J_m(A(a,Sa)).
 \tag{Double}\label{eq:sr-cycle-double}
\]
At zero, both sides are $F_m(F_m(0))$, because $A(0,1)=2$.
At a successor, KR applies $F_m$ twice to the left side, and JR plus
AD applies $F_m$ twice to the right side.  The induction equation
therefore supplies the step.

Now define the single finite SR term
\[
                              P(n):=K_n(n).
 \tag{P}\label{eq:sr-predecessor-term}
\]
Its zero equation is $P(0)=K_0(0)=0$.  For arbitrary $a$, specialize
the already proved identities at $m=Sa$ to obtain
\[
\begin{aligned}
 P(Sa)
 &=K_{Sa}(Sa)\\
 &=J_{Sa}(A(a,Sa)) &&\text{by Double}\\
 &=J_{Sa}(A(Sa,a)) &&\text{by commutativity of }A\\
 &=J_{Sa}(a)       &&\text{by Period}\\
 &=a              &&\text{by Run and }L(a,Sa)=1.
\end{aligned}
\]
Consequently
\[
                         \boxed{P(0)=0,\qquad P(Sa)=a.}
 \tag{Predecessor}\label{eq:sr-predecessor}
\]
The proof has not cancelled a successor or assumed an actual-equality
characteristic.  It works at arbitrary source elements, including
nonstandard elements of any model of the stated ordinary theory.
Successor injectivity can now be derived by applying $P$ to an equation
$Sa=Sb$.

\subsection{Ordinary initialization and both interpretation composites}

For any fixed finite update $h(u,\vec p)$, put
\[
 \widehat I(n,a;u.h)
 :=P\bigl(R(n,a;v.S(h[Pv/u]))\bigr),
 \tag{I}\label{eq:sr-initializer-compiler}
\]
with $v$ fresh.  Predecessor and SR give, by direct computation,
\[
 \widehat I(0,a;u.h)=a,\qquad
 \widehat I(Sn,a;u.h)=h[\widehat I(n,a;u.h)/u].
\]
In particular the ordinary selector is now available as
$C(n,a,b)=\widehat I(n,a;u.b)$.

The established initialized-scalar reconstruction then applies with
all its hypotheses satisfied: ordinary intuitionistic equality logic,
coverage and separation, literal-equation induction, successor, and
scalar iteration with arbitrary initial values.  It proves actual
equality reflection, compiles quantifier-free formulas into Boolean
terms, recovers contextual quantifier-free induction, and supplies
pairing and the compiler for every fixed primitive-recursive description.
Its written internal reconstruction, included elsewhere in the report,
is required for this conclusion; standard-model term completeness alone
would not suffice.  The predecessor step of that older reconstruction
is unnecessary here because Predecessor has already been proved.

Conversely, interpret each fixed SR description in ordinary iPRA by
primitive recursion with initial value $Sa$ and update $h$.  Then SR,
coverage, separation, and the translated literal-equation induction
instances are all iPRA theorems.  The language uses no function variables
or quantification over codes.

To check the stronger return condition, first anchor successor: the
zero instance $Sa=R(0,a;u.0)$ translates to ordinary $Sa$, and its
opposite translation is the source $Sa$. This uses the zero computation,
not uniqueness for an as-yet-unidentified recursor. Then compare one fixed description
with its double translation after earlier descriptions have been
identified.  For an SR description both sides have initial value $Sa$
and the same update $h$; EI on the literal equation comparing their
values proves their equality.  For an initialized scalar description,
the same argument uses its common initial value $a$.  The subsequent
primitive-recursion compiler has the same uniqueness argument for its
fixed descriptions.  In each argument the initializer and passive
parameters remain independent formal variables during induction; only
after the universally parameterized identity is proved are arbitrary
terms substituted for them.

Both interpretation domains are the entire numerical domain, equality
is actual equality, and the numerical composite isomorphisms have graph
$x=y$.  The ordinary logical translations are identities, so external
induction on formulas gives the full formula return translations.
Thus $T_{\rm SR}$ is equiconsistent and bi-interpretable with full
first-order iPRA.

\paragraph{What this proves about simplicity.}
The ordinary numerical grammar has just zero and the successor-seeded
scalar recursor. Section~\ref{sec:native-sr-bootstrap} additionally
translates this arithmetic proof into the native logical rules, before
ordinary connectives or an unshifted selector are available. The Boolean
successor decoder in that section removes the circularity that a direct
appeal to the established logical reconstruction would have introduced.
The numerical identity composites proved here therefore apply to the
direct native SR presentation as well.
